% Chapter 1 --- Introduction
% Drafted to match the voice and evidence of Chapters 3-5.

\section{Background and Motivation}

A private investor holding the S\&P~500 faces the same decision every
trading day: buy more, sell some, or do nothing. Most days the decision
barely matters. Some months it matters a great deal. In March 2025, the
worst month of the period studied in this dissertation, a passive investor
holding the S\&P~500 tracker fund SPY lost \$398.41 of a \$10,000 stake ---
nearly 4\% of their capital in a single month. An investor who had stepped
out of the market before that decline would have kept the money. An
investor who stepped out too often would have missed the gains of the other
23 months, which averaged over \$200 each. Timing the market well is
therefore a genuine sequential decision problem: every action changes the
position the next decision starts from, and the consequences of a decision
may only become visible weeks later.

Reinforcement learning (RL) is the branch of machine learning built for
exactly this structure. An RL agent acts, observes the outcome, and learns
from delayed rewards rather than from labelled examples. Trading appears to
be a natural fit: the market state is observable, actions are discrete
(buy, sell, hold), and the reward --- the change in portfolio wealth --- is
measurable to the cent. This apparent fit has produced a large and growing
body of work applying RL to trading.

There is, however, a catch that motivates the design of this entire study.
Equity markets rise more often than they fall. In a rising market, a
strategy that simply buys and holds earns a strong return while using no
market information at all. A learned agent that reports a positive return
has therefore proven very little. The meaningful questions are whether the
agent beats the simple strategy it competes against, and whether its
decisions actually depend on what it observes. Both questions must be
answered on data the agent has never seen, after realistic transaction
costs. This dissertation is built around answering them carefully.

\section{Problem Statement}

The core problem this dissertation addresses is the following: can standard
reinforcement learning algorithms learn a trading policy for a single
liquid asset that outperforms simple benchmark strategies on unseen data,
once transaction costs are accounted for?

Behind this headline question sits a methodological problem that is just as
important. Financial return alone cannot show whether an agent has learned
anything. A policy that buys whenever buying is allowed performs well in a
rising market without reading a single market feature. Evaluating learned
traders therefore requires a second, behavioural axis: determining whether
a policy's actions change in response to its observed state at all. Without
that check, a strong return from an unconditional buyer could be mistaken
for evidence of learned skill. This dissertation develops and applies such
a check alongside the financial comparison.

\section{Research Questions and Objectives}

The aim of this dissertation is to evaluate, fairly and reproducibly,
whether three standard reinforcement learning algorithms --- REINFORCE,
Deep Q-learning (DQN) and Proximal Policy Optimisation (PPO) --- can learn
profitable and genuinely state-dependent trading policies for the SPY
exchange-traded fund. The evaluation is structured around four research
questions, stated here and answered in Chapter 5:

\begin{enumerate}
\def\labelenumi{\arabic{enumi}.}
\item
  \textbf{Performance:} do the learned agents outperform appropriate
  simple benchmarks on unseen data after accounting for transaction
  costs?
\item
  \textbf{Behaviour:} do the learned agents make use of the market and
  portfolio information provided in the state representation?
\item
  \textbf{State representation:} does adding a drawdown feature to the
  main nine-feature state affect the learned behaviour or financial
  performance?
\item
  \textbf{Reward:} does a risk-aware reward improve the balance between
  return and drawdown, rather than simply encouraging the agent to trade
  less?
\end{enumerate}

These questions translate into seven measurable objectives, each traceable
to a specific part of this dissertation:

\begin{enumerate}
\def\labelenumi{O\arabic{enumi}.}
\item
  Build a reproducible daily trading environment for SPY with monthly
  episodes, realistic transaction fees and action masking, together with
  an automated verification suite (Chapter 4).
\item
  Implement the three learning algorithms and verify each implementation
  before any trading experiment (Chapters 3 and 4).
\item
  Calibrate the fixed trade size on validation data so that the agents
  can reach at least 90\% of the buy-and-hold benchmark's market
  exposure, making the comparison fair (Section 5.3).
\item
  Determine, with a behavioural diagnostic, whether each learned policy
  actually uses its state observations (Section 5.4).
\item
  Measure the effect of adding the drawdown feature to the state while
  holding everything else fixed (Section 5.5).
\item
  Compare the final agents against the benchmarks on 24 held-out months,
  using six random seeds per configuration (Section 5.6).
\item
  Test the sensitivity of the findings to the reward formulation and to
  the level of transaction costs (Sections 5.7 and 5.8).
\end{enumerate}

\section{Scope}

The study is deliberately narrow so that each comparison within it can be
clean. All experiments use one asset: SPY, the exchange-traded fund
tracking the S\&P~500 index, observed at daily resolution from 2018 to
2025 and split chronologically into training (2018--2022), validation
(2023) and held-out test data (2024--2025). Each monthly episode starts
with \$10,000 of capital, and every trade pays a 0.05\% transaction fee
unless stated otherwise. Actions are discrete --- buy a fixed fraction of
the initial capital, sell that fraction, or hold --- and an action mask
prevents invalid trades such as selling with no position. Three algorithms
are compared: REINFORCE and Deep Q-learning implemented from scratch, and
PPO from an established library as a comparator.

Multi-asset portfolios, continuous position sizing, intraday data and live
deployment are out of scope. These boundaries keep the number of
experimental variables small enough that every observed difference can be
attributed to a deliberate change of exactly one factor.

\section{Contributions}

This dissertation makes four contributions:

\begin{enumerate}
\def\labelenumi{\arabic{enumi}.}
\item
  \textbf{A verified evaluation framework for learned trading agents.}
  Every design decision is made on validation data and frozen before the
  held-out test months are touched. The environment, features, rewards
  and benchmarks are covered by an automated test suite, and the
  corrections this suite forced --- including a benchmark implementation
  error worth \$38.95 per month --- are documented rather than hidden
  (Chapter 4).
\item
  \textbf{A behavioural state-dependence diagnostic.} The diagnostic
  separates policies that respond to market information from policies
  that merely follow the pattern of legal actions. It reveals a clean
  split that return figures alone would hide: all six Deep Q-learning
  seeds used their observations, while none of the twelve policy-gradient
  seeds did (Section 5.4).
\item
  \textbf{A fairness calibration for the action model.} The fixed trade
  size is shown to be part of the problem definition rather than an
  implementation detail: an uncalibrated size caps the agent's achievable
  exposure at 0.690 against the benchmark's 0.910 and would have made the
  main comparison unfair before any learning took place (Section 5.3).
\item
  \textbf{A carefully validated negative result with a mechanistic
  explanation.} None of the three agents beat buy-and-hold on the 24
  held-out months, and the behavioural analysis explains why: the
  policies that earned the most did so by ignoring their observations,
  while the one method that used its observations reduced participation
  rather than timing the market. The result is negative but the evidence
  base --- six seeds per configuration, behavioural diagnostics, and
  sensitivity checks over states, rewards and fees --- makes it a finding
  rather than a failure (Chapters 5 and 6).
\end{enumerate}

\section{Dissertation Structure}

The remainder of this dissertation is organised as follows. Chapter 2
reviews the literature on reinforcement learning for trading and positions
this study's evaluation-first approach within it. Chapter 3 presents the
mathematical formulation: the trading environment as a Markov decision
process, the state and reward designs, and the three learning algorithms.
Chapter 4 describes the system implementation, the data pipeline, the
verification suite and the corrections it forced. Chapter 5 reports the
experiments and analyses the results around the four research questions.
Chapter 6 concludes with the main findings, a critical appraisal of the
work, its limitations, and directions for future work.
